\section{USAMO 1973}
        \begin{enumerate}
        	\item (\textit{USAMO 1973}) Two points $P$ and $Q$ lie in the interior of a regular tetrahedron $ABCD$. Prove that angle $PAQ<60^o$.


 



\textbf{Solution 1}

Let the side length of the regular tetrahedron be $a$. Link and extend $AP$ to meet the plane containing triangle $BCD$ at $E$; link $AQ$ and extend it to meet the same plane at $F$. We know that $E$ and $F$ are inside triangle $BCD$ and that $\angle PAQ = \angle EAF$

 Now let’s look at the plane containing triangle $BCD$ with points $E$ and $F$ inside the triangle. Link and extend $EF$ on both sides to meet the sides of the triangle $BCD$ at $I$ and $J$, $I$ on $BC$ and $J$ on $DC$. We have $\angle EAF < \angle IAJ$

 But since $E$ and $F$ are interior of the tetrahedron, points $I$ and $J$ cannot be both at the vertices and $IJ < a$, $\angle IAJ < \angle BAD = 60$. Therefore, $\angle PAQ < 60$.

 Solution with graphs posted at

 \href{https://artofproblemsolving.comhttp://www.cut-the-knot.org/wiki-math/index.php?n=MathematicalOlympiads.USA1973Problem1}{http://www.cut-the-knot.org/wiki-math/index.php?n=MathematicalOlympiads.USA1973Problem1}

 



\textbf{Solution 2 (similar to solution 1, but full proof)}

Let the side length of the regular tetrahedron $ABCD$ be $s$.

 Extend $AP$ and $AQ$ to intersect the plane $BCD$ at points $P'$ and $Q'$ inside $\Delta BCD$, respectively. We can observe that $\angle PAQ=\angle P'AQ'$.

 Extend $P'Q'$ to intersect the sides of $\Delta BCD$ at $P''$ and $Q''$. We can observe that $\angle P'AQ'<\angle P''AQ''$. We assume that, WLOG, $P''$ is on $BD$ and $Q''$ is on $BC$.

 Let the ratios $\frac{BP''}{BD}=\mu$ and $\frac{BQ''}{BC}=\lambda$. We can see that $\mu,\lambda\in (0,1)$. It follows from Stewart's theorem that$AP''^2=AD^2\mu+AB^2(1-\mu)-\mu(1-\mu)BD^2=(\mu^2-\mu+1)s\\ AQ''^2=AC^2\lambda+AB^2(1-\lambda)-\lambda(1-\lambda)BC^2=(\lambda^2-\lambda^2+1)s$

 Furthermore, we have$P''C^2=CD^2\mu+CB^2(1-\mu)-\mu(1-\mu)BD^2=(\mu^2-\mu+1)s$

 Using this, we can conclude that$P''Q''^2=P''C^2\lambda+(BD\mu)^2(1-\lambda)-\lambda(1-\lambda)BC^2\\ =(\mu\lambda-\mu-\lambda+2)s^2$

 By cosine law, we have$\cos \angle P''AQ''=\frac{P''A^2+Q''A^2-P''Q''^2}{2P''A\cdot Q''A}\\ =\frac{\mu\lambda-\mu-\lambda+2}{2\sqrt{(\mu^2-\mu+1)(\lambda^2-\lambda+1)}}$for $\mu,\lambda\in(0,1)$.

 We first find a lower bound of the numerator. We write $\mu\lambda-\mu-\lambda+2=(\lambda-1)\mu-\lambda+1$ as a linear function of $\mu$ by letting $\lambda$ be constant. Since $\lambda\in (0,1)$, the coefficient $\lambda-1$ of $\mu$ in this linear function is negative, meaning this expression is least when $\mu$ is least "at" $1$. Hence, $\mu\lambda-\mu-\lambda+2<1$.

 Now we find an upper bound of the denominator. We complete the square in the expression $2\sqrt{(\mu^2-\mu+1)(\lambda^2-\lambda+1)}=2\sqrt{((\mu-\frac{1}{2})^2+\frac{3}{4})((\lambda-\frac{1}{2})^2+\frac{3}{4})}$ which is least when $\lambda$ and $\mu$ "are" both 1. Hence, $2\sqrt{(\mu^2-\mu+1)(\lambda^2-\lambda+1)}>2$.

 Hence, $\cos \angle P''AQ''<\frac{1}{2}$, implying that $\angle P''AQ''<60^\circ$.

 Therefore, $\angle PAQ=\angle P'AQ'<\angle P''AQ''<60^\circ$.

 <i>Alternate solutions are always welcome. If you have a different, elegant solution to this problem, please \href{https://artofproblemsolving.com/wiki/index.php?title=1973\_USAMO\_Problems/Problem\_1&action=edit}{add it} to this page.</i>

 hurdler: Remark on solution 1: This proof is not rigorous, in the very last step. The last step needs more justification.

 


	\item (\textit{USAMO 1973}) Let $\{X_n\}$ and $\{Y_n\}$ denote two sequences of integers defined as follows:


  $X_0=1$, $X_1=1$, $X_{n+1}=X_n+2X_{n-1}$ $(n=1,2,3,\dots),$

 $Y_0=1$, $Y_1=7$, $Y_{n+1}=2Y_n+3Y_{n-1}$ $(n=1,2,3,\dots)$.

 Thus, the first few terms of the sequences are:


  $X:1, 1, 3, 5, 11, 21, \dots$,

  $Y:1, 7, 17, 55, 161, 487, \dots$.

 Prove that, except for the "1", there is no term which occurs in both sequences.


 



\textbf{Solution}

We can look at each sequence $\bmod{8}$:

 $X$: $1$, $1$, $3$, $5$, $3$, $5$, $\dots$,

 $Y$: $1$, $7$, $1$, $7$, $1$, $7$, $\dots$.

 
                \begin{itemize}
                    \item Proof that $X$ repeats $\bmod{8}$:
                \end{itemize}


                The third and fourth terms are $3$ and $5$ $\bmod{8}$. Plugging into the formula, we see that the next term is $11\equiv 3\bmod{8}$, and plugging $5$ and $3$, we get that the next term is $13\equiv 5\bmod{8}$. Thus the sequence $X$ repeats, and the pattern is $3,5,3,5,\dots$.

 
                \begin{itemize}
                    \item Proof that $Y$ repeats $\bmod{8}$:
                \end{itemize}


                The first and second terms are $1$ and $7$ $\bmod{8}$. Plugging into the formula, we see that the next term is $17\equiv 1\bmod{8}$, and plugging $7$ and $1$, we get that the next term is $23\equiv 7\bmod{8}$. Thus the sequence $Y$ repeats, and the pattern is $1,7,1,7,\dots$.

 


 Combining both results, we see that $X_i$ and $Y_j$ are not congruent $\bmod{8}$ when $i\geq 3$ and $j\geq 2$. Thus after the "1", the terms of each sequence are not equal.

 


 



\textbf{Solution 2}

We can solve this problem by finding a particular solution for each linear recurrence.

 $X_{n+1} - X_n - 2X_{n-1} = 0$ with characteristic polynomial

 $X^2 - X - 2 = (X - 2)(X + 1)$, so $X = -1, 2$

 $X_n = A(2)^n + B(-1)^n$

 After plugging in to find the particular solution:

 $X_0 = 1 = A + B$

 $X_1 = 1 = 2A - B$

 $A = \frac{1}{3}$ and $B = \frac{2}{3}$, so

 $X_n = \frac{1}{3}(2)^n + \frac{2}{3}(-1)^n$

 Doing the same for $Y_n$, we get

 $Y_n = 2(3)^n - (-1)^n$

 We know they're equal at $n = 0$, so let's set them equal and compare.

 $2(3)^n - (-1)^n = \frac{1}{3}((2)^n + 2(-1)^n)$

 $6(3)^n - 3(-1)^n = (2)^n + 2(-1)^n$

 $6(3)^n - 2^n = 5(-1)^n$

 By induction, we know in general that $(\alpha + 1)^n > \alpha^n$ for all $\alpha, n > 0$, so $6(3)^n > 2^n$ for all $n > 1$. 

 Therefore, the left hand side is always increasing and will never equal 5 again, so equality between $X_n$ and $Y_n$ only holds when $n = 0$, so they only share 1 term.

 


 



\textbf{Solution 3 (simple)}

First, $X_n$ and $Y_n$ are both positive for all nonnegative integers $n$. This is since $X_1$, $X_2$, $Y_1$, and $Y_2$ are all positive, and if $X_n$ and $X_{n-1}$ is positive then $X_{n+1}$ as a sum of two positive numbers is positive, and similarly for $Y_n$. Simple.

 Next, consider $Z_n = Y_n - X_n$. We have $Z_{n+1} = Y_n + Z_n + 3 Z_{n-1} + X_{n-1}$. Another easy inductive argument like above shows that $Z_n >0$ if $n>0$. But if $Z_n >0$ then $Z_n \ne 0$, so $X_n \ne Y_n$ as desired.

 
<i>Alternate solutions are always welcome. If you have a different, elegant solution to this problem, please \href{https://artofproblemsolving.com/wiki/index.php?title=1973\_USAMO\_Problems/Problem\_2&action=edit}{add it} to this page.</i>

 


	\item (\textit{USAMO 1973}) Three distinct vertices are chosen at random from the vertices of a given regular polygon of $(2n+1)$ sides. If all such choices are equally likely, what is the probability that the center of the given polygon lies in the interior of the triangle determined by the three chosen random points?


 



\textbf{Solution}

There are $\binom{2n+1}{3}$ ways how to pick the three vertices. We will now count the ways where the interior does NOT contain the center. These are obviously exactly the ways where all three picked vertices lie among some $n+1$ consecutive vertices of the polygon.We will count these as follows: We will go clockwise around the polygon. We can pick the first vertex arbitrarily ($2n+1$ possibilities). Once we pick it, we have to pick $2$ out of the next $n$ vertices ($\binom{n}{2}$ possibilities).

 Then the probability that our triangle does NOT contain the center is 
 \[p = \frac{ (2n+1){\binom{n}{2}} }{ {\binom{2n+1}{3} } } = \frac{ (1/2)(2n+1)(n)(n-1) }{ (1/6)(2n+1)(2n)(2n-1) } =  \frac{ 3n-3 }{ 4n-2 }\]

 And then the probability we seek is 
 \[1-p =  \frac{ (4n-2) - (3n-3) }{ (4n-2) } = \boxed{\frac{n+1}{4n-2}}\]

 \textbf{Perplexing Edge Case} However, for $n = 1$, (a rare edge case), the circumcenter of the triangle must lie in the center of the triangle, which is notably difficult to prove if not assumed, resulting in a probability of $1$.

 Observe that the probability tends to $\frac{1}{4}$ as $n$ approaches infinity. Why does this occur?

 
<i>Alternate solutions are always welcome. If you have a different, elegant solution to this problem, please \href{https://artofproblemsolving.com/wiki/index.php?title=1973\_USAMO\_Problems/Problem\_3&action=edit}{add it} to this page.</i>

 


	\item (\textit{USAMO 1973}) Determine all the \href{https://artofproblemsolving.com/wiki/index.php?title=Root}{roots}, \href{https://artofproblemsolving.com/wiki/index.php?title=Real}{real} or \href{https://artofproblemsolving.com/wiki/index.php?title=Complex}{complex}, of the system of simultaneous \href{https://artofproblemsolving.com/wiki/index.php?title=Equation}{equations}


 $x+y+z=3$,
$x^2+y^2+z^2=3$,


 
$x^3+y^3+z^3=3$.

 



\textbf{Solution 1}

Let $x$, $y$, and $z$ be the \href{https://artofproblemsolving.com/wiki/index.php?title=Root}{roots} of the \href{https://artofproblemsolving.com/wiki/index.php?title=Cubic\_polynomial}{cubic polynomial} $t^3+at^2+bt+c$. Let $S_1=x+y+z=3$, $S_2=x^2+y^2+z^2=3$, and $S_3=x^3+y^3+z^3=3$. From this, $S_1+a=0$, $S_2+aS_1+2b=0$, and $S_3+aS_2+bS_1+3c=0$. Solving each of these, $a=-3$, $b=3$, and $c=-1$. Thus $x$, $y$, and $z$ are the roots of the polynomial $t^3-3t^2+3t-1=(t-1)^3$. Thus $x=y=z=1$, and there are no other solutions.

 



\textbf{Solution 2}

Let $P(t)=t^3-at^2+bt-c$ have roots x, y, and z. Then 
 \[0=P(x)+P(y)+P(z)=3-3a+3b-3c\] using our system of equations, so $P(1)=0$. Thus, at least one of x, y, and z is equal to 1; without loss of generality, let $x=1$. Then we can use the system of equations to find that $y=z=1$ as well, and so $\boxed{(1,1,1)}$ is the only solution to the system of equations.

 



\textbf{Solution 3}

Let $a=x-1,$ $b=y-1$ and $c=z-1.$ Then 
 \[a+b+c=0,\]
 \[a^2+b^2+c^2=0,\]
 \[a^3+b^3+c^3=0.\]We have
 \begin{align*} 0&=(a+b+c)^3\\ &=(a^3+b^3+c^3)+3a^2(b+c)+3b^2(a+c)+3c^2(a+b)+6abc\\ &=0-3a^3-3b^3-3c^3+6abc\\ &=6abc. \end{align*}
Then one of $a, b$ and $c$ has to be 0, and easy to prove the other two are also 0. So $\boxed{(1,1,1)}$ is the only solution to the system of equations.

 J.Z.

 



\textbf{Solution 4}

We are going to use Intermediate Algebra Techniques to solve this equation.

 Let's start with the first one: $x+y+z=3$. This will be referred as the FIRST equation.

 We are going to use the first equation to relate to the SECOND one ($x^2+y^2+z^2=3$) and the THIRD one ($x^3+y^3+z^3=3)$.

 Squaring this equation: $x^2+y^2+z^2+2xy+2yz+2xz=9$

 Subtracting this equation from the 2nd equation in the problem, we have $2xy+2yz+2xz=6$, so $xy+xz+yz=3$.

 Now we try the same idea with the cubed terms. Cube the first equation:

 $x^3+y^3+z^3+3xy^2+3xz^2+3yx^2+3yz^2+3zx^2+3zy^2=27$. Plug in $x^3+y^3+z^3=3$ and factor partially:

 $3+3(x^2(y+z)+y^2(x+z)+z^2(x+y))+6xyz=27$.

 Now here is the key step. Note that $z=3-x-y, y=3-x-z, x=3-y-z$. So we are going to substitute $y+z, x+z, x+y$ for each of the expressions and we get: $-x^3+3x^2-y^3+3y^2-z^3+3z^2=8-2xyz$ (I rearranged it a bit).

 Resubstituting in the second and third equation: $-3+3(3)=8-2xyz$. So $xyz=1$.

 So now we have three equations for the elementary symmetric sums of $x,y,z$:

 Equation 4: $x+y+z=3$ (this is also equation 1)

 Equation 5: $xy+yz+xz=3$

 Equation 6: $xyz=1$.

 If we call the solutions of $t^3-3t^2+3t-1=0$ (Equation 7) $a,b,c$, then $x,y,z$ are the three roots $a,b,c$ but in some order. Notice that Equation 7 can be factored as $(t-1)^3=0$, which means that $t=1$. Therefore $(x,y,z)$ are permutations of $(1,1,1)$ in some order, which can be only $(1,1,1)$. (This step uses Vieta's formulas)

 Therefore, the only solution is $\boxed{(x,y,z)=(1,1,1)}$.

 



\textbf{Solution 5}

From the question

 Equation 1 : $x+y+z=3$

 Equation 2 : $x^2+y^2+z^2=3$

 Equation 3 : $x^3+y^3+z^3=3$

 From Eq(1) and Eq(2),

 Equation 4 : $(x+y+z)^2=9 \implies x^2+y^2+z^2+2\sum xy = 9 \implies \sum xy = 3$

 From Eq(1), Eq(3), and the fact that $(x+y+z)^3=x^3+y^3+z^3+3\prod(x+y)$,

 Equation 5 : $(x+y+z)^3=x^3+y^3+z^3+3\prod(x+y)= 27 \implies \prod(x+y)=8$

 As from Eq(1),  $\ x+y=3-z$,$\ y+z=3-x$,  $\ z+x = 3- y$

 Equation 6 : $\prod(3-x) =8$

 Now let there be a cubic function $f(t)$ with the roots $x,y,z$ so from the relation between roots and coefficients,

 Equation 7 : $f(t)= t^3-3t^2+3t-d$   (As we don't know the product let the constant be $d$)

 As $f(t) = (t-x)(t-y)(t-z)$ then at $t=3$,

 Equation 8 : $f(3)=9-d=(3-x)(3-y)(3-z)$

 And from Eq(6),

 Equation 9 : $d= -1$

 So finally,

 Equation 10 : $f(t)=t^3-3t^2+3t-1=(t-1)^{3}$

 So, $\ (x,y,z)=(1,1,1)$

 <i>Alternate solutions are always welcome. If you have a different, elegant solution to this problem, please \href{https://artofproblemsolving.com/wiki/index.php?title=1973\_USAMO\_Problems/Problem\_4&action=edit}{add it} to this page.</i>

 


	\item (\textit{USAMO 1973}) Show that the cube roots of three distinct prime numbers cannot be three terms (not necessarily consecutive) of an arithmetic progression.


 



\textbf{Solution}

Assume that the cube roots of three distinct prime numbers can be three terms of an arithmetic progression. Let the three distinct prime numbers be $p$, $q$, and $r.$ WLOG, let $p<q<r.$

 Then, 

 
 \[q^{\dfrac{1}{3}}=p^{\dfrac{1}{3}}+md\]

 
 \[r^{\dfrac{1}{3}}=p^{\dfrac{1}{3}}+nd\]

 where $m$, $n$ are distinct integers, and $d$ is the common difference in the progression. Then we have

 
 \[\begin{array}{rcl} nq^{\dfrac{1}{3}}-mr^{\dfrac{1}{3}}&=&(n-m)p^{\dfrac{1}{3}} \\  \\n^{3}q-3n^{2}mq^{\dfrac{2}{3}}r^{\dfrac{1}{3}}+3nm^{2}q^{\dfrac{1}{3}}r^{\dfrac{2}{3}}-m^{3}r&=&(n-m)^{3}p\\  \\3nm^{2}q^{\dfrac{1}{3}}r^{\dfrac{2}{3}}-3n^{2}mq^{\dfrac{2}{3}}r^{\dfrac{1}{3}}&=&(n-m)^{3}p+m^{3}r-n^{3}q\\  \\(3nmq^{\dfrac{1}{3}}r^{\dfrac{1}{3}})(mr^{\dfrac{1}{3}}-nq^{\dfrac{1}{3}})&=&(n-m)^{3}p+m^{3}r-n^{3}q\\  \\(3nmq^{\dfrac{1}{3}}r^{\dfrac{1}{3}})((m-n)p^{\dfrac{1}{3}})&=&(n-m)^{3}p+m^{3}r-n^{3}q\\  \\q^{\dfrac{1}{3}}r^{\dfrac{1}{3}}p^{\dfrac{1}{3}}&=&\dfrac{(n-m)^{3}p+m^{3}r-n^{3}q}{(3mn)(m-n)}\\  \\(pqr)^{\dfrac{1}{3}}&=&\dfrac{(n-m)^{3}p+m^{3}r-n^{3}q}{(3mn)(m-n)} \end{array}\]

 Because $p$, $q$, $r$ are distinct primes, $pqr$ is not a perfect cube. Thus, the LHS is irrational but the RHS is rational, which is a contradiction. So, the cube roots of three distinct prime numbers cannot be three terms of an arithmetic progression.

 



\textbf{Solution 2}

Assume the contrary and suppose that the cube roots of three distinct prime numbers, $p_1^\frac{1}{3}$, $p_2^\frac{1}{3}$, and $p_3^\frac{1}{3}$, form an increasing arithmetic progression. Then we can write $p_2^\frac{1}{3} = p_1^\frac{1}{3} + md$ and $p_3^\frac{1}{3} = p_1^\frac{1}{3} + nd$ for some positive integral $d, m, n$ with $m < n$. Thus,
 \[\frac{p_3^\frac{1}{3} - p_1^\frac{1}{3}}{p_2^\frac{1}{3} - p_1^\frac{1}{3}} = \frac{md}{nd} = \frac{m}{n} = r,\]for some real $r$. Because $m$ and $n$ are integral, r is rational.

 This equation rearranges to 
 \[p_3^\frac{1}{3} = rp_2^\frac{1}{3} - (r-1)p_1^\frac{1}{3}.\]

 Note that if $a-b=x$, then 
 \[(a-b)^3 = a^3 - b^3 - 3ab(a-b) = a^3 - b^3 - 3abx\] from the binomial theorem. As a result, cubing both sides of the equation gives
 \[p_3 = r^3p_2 - (r-1)^3p_1 - 3r(r-1)(p_1p_2p_3)^\frac{1}{3}.     (*)\]

 \textit{Lemma:} The sum, product, and difference of two rational numbers are all rational, and that of a rational and irrational number are all irrational.

 Proof: The first part is trivial if we write the fractions as $\frac{m}{n}$ and $\frac{p}{q}$, and find the common denominators or simplify fractions asneeded. The second part follows by contradiction: if the sum of a rational number and irrational number is rational, then the irrational number must be the difference of two rational numbers, which is rational. This is a contradiction; likewise for difference and product.

 The LHS \textit{(left-hand side)} of (*) is clearly rational (in fact, integral), and so the RHS must also be rational. This is the case for the first two terms $r^3p_2$ and $-(r-1)^3p_1$. But it is not the case for the third term. The number $(p_1p_2p_3)^\frac{1}{3}$ is irrational because $p_1p_2p_3$, the product of three distinct primes, cannot be a perfect cube. Hence, $-3r(r-1)(p_1p_2p_3)^\frac{1}{3}$, the product of a rational and an irrational number, is irrational, and the RHS is therefore the sum of two rational terms and one irrational term. This is irrational by the lemma, and hence we have a rational equalling an irrational, contradiction. It follows that the cube roots of three distinct prime numbers cannot form an arithmetic sequence.

 


\end{enumerate}